$A(U \wedge I)$ is a bivector of $V_{\C}$, so by Proposition \claimref{standard-model.module-factorises-even-ground-elements-act}(1) it acts as $A_{W} \otimes 1$. For a product form and states $e_{a} \otimes \psi$, $e_{b} \otimes \varphi$ with $e_{a}, e_{b}$ in rung sectors $j, j'$: $\ip{e_{a} \otimes \psi}{(A_{W} \otimes 1)(e_{b} \otimes \varphi)} = \ip{e_{a}}{e_{b}}_{R}\ip{\psi}{A_{W}\varphi}_{W}$, which vanishes for $j \ne j'$ (different rung numbers are orthogonal under the rung Fock form) and for $a \ne b$ within a sector (orthonormality). Expanding $\Psi$ in the basis $e_{a} \otimes \psi_{a}$ gives (i) and (ii) for the moment; the grip is linear in the moment through the nondegenerate pairing, so the sums pass to $P$; and $U(2)$ acts as $1 \otimes \Lambda^{\bullet}(A)$, which commutes with $A_{W} \otimes 1$ and is unitary for the rung form, so the moment is invariant. (iii) restates (ii): the right-hand side does not mention $a$ beyond the index of summation, and $\binom{2}{j}$ is the only sector datum that enters. The script checks (i) on three Gaussian-rational states of the module and eight bivectors, and (ii) on the doublet sector and on both singlet sectors, in the Fock form; the proof is the same one line for any product form.
